\usepackage[paperwidth=539.5bp,paperheight=711.85bp,left=49bp,right=49bp,top=43bp,bottom=55bp,footskip=27bp]{geometry}
\usepackage{amsmath,amssymb,bm,mathtools}
\usepackage{graphicx,adjustbox,xcolor}
\usepackage{fancyhdr}
\usepackage[normalem]{ulem}
\usepackage{array,booktabs,multirow}
\usepackage{fontspec}
\setmainfont{lmroman10-regular.otf}[BoldFont=lmroman10-bold.otf,ItalicFont=lmroman10-italic.otf,BoldItalicFont=lmroman10-bolditalic.otf]
\setsansfont{lmsans10-regular.otf}[BoldFont=lmsans10-bold.otf]
\setmonofont{lmmono10-regular.otf}
\newfontfamily\cyrillicfont{DejaVuSerif.ttf}[Path=fonts/,BoldFont=DejaVuSerif-Bold.ttf,ItalicFont=DejaVuSerif-Italic.ttf,BoldItalicFont=DejaVuSerif-BoldItalic.ttf]
\newcommand{\textcyrillic}[1]{{\cyrillicfont #1}}
\setlength{\parindent}{1.8em}
\setlength{\parskip}{5pt plus 1pt}
\setlength{\abovedisplayskip}{7pt}
\setlength{\belowdisplayskip}{7pt}
\setlength{\emergencystretch}{2em}
\fancyhf{}
\renewcommand{\headrulewidth}{0pt}
\renewcommand{\footrulewidth}{0pt}
\newcommand{\CurrentPrintedPage}{}
\fancypagestyle{scan}{\fancyhf{}\fancyfoot[L]{\ttfamily\fontsize{8}{10}\selectfont FTD-HC-23-259-71}\fancyfoot[C]{\CurrentPrintedPage}}
\pagestyle{scan}
\newsavebox{\ScanContent}
\NewDocumentEnvironment{scanpage}{m m +b}{%
  \clearpage\begingroup
  \gdef\CurrentPrintedPage{#2}%
  \sbox{\ScanContent}{\begin{minipage}[t]{\textwidth}%
    \fontsize{10.5}{14.4}\selectfont
    #3\par
  \end{minipage}}%
  \typeout{SCANPAGE #1 LABEL #2 HEIGHT \the\dimexpr\ht\ScanContent+\dp\ScanContent\relax LIMIT \the\textheight}%
  \noindent\begin{adjustbox}{max totalsize={\textwidth}{\textheight},valign=t}\usebox{\ScanContent}\end{adjustbox}%
  \par\clearpage\endgroup
}{}
\newcommand{\chaptertitle}[2]{\begin{center}\vspace*{12pt}CHAPTER #1\par\vspace{15pt}\textbf{#2}\par\vspace{15pt}\end{center}}
\newcommand{\sectiontitle}[2]{\par\medskip\begin{center}\S\ #1.\quad\uline{#2}\end{center}\medskip}
\newcommand{\unclear}[1]{\textcolor{orange!75!black}{#1}}
\DeclareMathOperator{\tg}{tg}
\DeclareMathOperator{\ctg}{ctg}
\DeclareMathOperator{\rot}{rot}
\DeclareMathOperator{\Div}{div}
\DeclareMathOperator{\grad}{grad}
\DeclareMathOperator{\sgn}{sgn}
\DeclareMathOperator{\RePart}{Re}
\DeclareMathOperator{\ImPart}{Im}
